\section{Applicability and Conditions}
\label{sec:limits}
\sys{} does not apply to every EP implementation, and being precise about where it
applies is part of the contribution. This section states the conditions the design
requires, then classifies the baseline paths against them.

\subsection{Four conditions for correctness}
\label{sec:conditions}
\paragraph{Generation-scoped resources.}
Every \data{} unit must land in staging that is valid for the current generation, and a
slot becomes reusable only after its last consumer finishes reading. Emitting \done{}
does not release the buffer: the destination GPU's successor kernels may still be
reading it, so notification and recycling are separate events with separate timestamps.

\paragraph{Ordered visibility.}
Observing \done{} must imply that all data it covers is visible to the consumer
(Equation~\eqref{eq:visible}). This is a delivery-contract requirement, not a
consequence of transmission order. Ordering is generally preserved across nodes but not
within a node, so a fence is needed wherever the covered data can be reordered relative
to the notification.

\paragraph{All managed traffic traverses the observer.}
The node can only reason about traffic it sees. A bypass path that skips it would make
the expected-count condition unsound, so we scope the design to configurations where the
relevant traffic is managed. Multiple rails or switches require partitioned completion
domains and an explicit join.

\paragraph{Bounded capacity with control progress.}
Staging is finite, and pressure propagates back along the same path. Unadmitted traffic
waits at the sender rather than being dropped, so exhausted capacity degrades coverage
rather than correctness. Critically, \ready{} messages must not be blocked behind queued
\data{}: a full staging pool that stalls the very messages needed to drain it would
deadlock. Reserved control resources or an equivalent separation are therefore required,
particularly under link-level backpressure~\cite{pfc}.

\begin{table*}[t]
\centering\small
\renewcommand{\arraystretch}{1.15}
\begin{tabularx}{\textwidth}{@{}>{\raggedright\arraybackslash}p{3.2cm}YYY@{}}
\toprule
Path & Before \data{} & Tail behavior & Role in this paper\\
\midrule
Cross-node Direct\newline Dispatch and Combine & EP-wide readiness barrier & Sender completion and local join precede the EP-wide signal exchange & Primary scenario for both mechanisms; tail saving depends on the exposed interval.\\
Hybrid Dispatch & Concurrent SU and SO readiness barriers & SU completion after forwarding; no additional SO barrier & Secondary scenario for readiness overlap.\\
Hybrid Combine & Concurrent SU and SO readiness barriers & Per-channel flush followed by SO completion signaling & Related completion path; assessed separately from Direct.\\
NCCL EP HT Dispatch & Routing bitmap AllGather and metadata scan & Chunk notifications, SU completion, and a cross-round guard & Metadata-dependent baseline; the guard protects buffer reuse.\\
NCCL EP LL Dispatch & No independent readiness barrier & Count notification follows local PUT issuance & Outside the primary two-boundary scenario.\\
Single-node Direct & NVLink readiness barrier & NVLink completion barrier & Scale-up reference; outside the managed Ethernet path.\\
\bottomrule
\end{tabularx}
\caption{Synchronization boundaries of the baseline communication paths~\cite{deepepcode,ncclhtcode,ncclllcode}. The applicable mechanism depends on the path's entry and completion dependencies. SU/SO denote scale-up/scale-out.}
\label{tab:paths}
\end{table*}

\subsection{Completion and visibility}
\label{sec:tailaudit}
Direct's sender completion, local coordination, and signal exchange provide the baseline for output-generated \done{}~\cite{gdakicode}. Section~\ref{sec:directtail} quantifies the exposed delay using data visibility and signal observation as distinct events.

Hybrid Dispatch ends with a scale-up barrier because its forwarders have already consumed the scale-out tokens~\cite{deepepcode}. NCCL HT publishes chunk signals on its data context and protects cross-round reuse with a guard; LL sends its count signal after local PUT issuance~\cite{ncclhtcode,ncclllcode}. These completion structures explain the scope of Table~\ref{tab:paths}.

\paragraph{Local bypass and hierarchical completion.}
The consumer joins network \done{} with required local NVLink writes and layout readiness; Hybrid also requires rail-to-target forwarding. These unobserved local paths retain their completion and buffer-ownership checks. If they dominate, earlier network notification does not advance successor release.

\paragraph{Performance and implementation boundary.}
Satisfying the four conditions makes the design correct, not profitable. Preallocation, staging bandwidth, output service, and notification overhead determine the net gain, and they are exactly the terms $\epsilon$ collects in \secref{sec:sensitivity}. Independent source-lifetime tracking must preserve safe reuse once sender completion is removed from the target-notification path, which is an added obligation rather than a relaxed one. Staging pressure and retained local dependencies bound the overlap that remains available.
